\section{Algebraic complexity and exact certificates}\label{s2:algebraic}

The quadratic equations of pooling give an existential-real upper bound,
but their network structure might conceivably make exact decision easier.
We show that it does not: even one attribute, binary source qualities,
and small local degrees suffice for $\ETR$-completeness. A different
construction puts the entire nonlinear interaction in one pool, using
bypasses and a growing number of attributes. We then identify why these
results coexist with $\NP$ certificates when suitable dimensions are
fixed. All results concern the standard, finite-capacity model of
Section~\ref{s1:model}.

The distinction from the earlier complexity results is substantive.
\citet[Theorems~3--5]{haugland2016-the-computational-complexity-of-the}
establishes $\NP$-hardness with one upper-bounded attribute under several
restrictions on inputs and products. Those reductions establish
combinatorial hardness. Our reductions instead realize the arithmetic
operations of an $\ETR$-complete problem, using exact product quality
specifications. To the best of our knowledge, the pooling literature
reviewed here contains no previous $\ETR$-completeness result for either
the bounded-degree, one-attribute class of Theorem~\ref{s2:bounded} or
the one-pool class of Theorem~\ref{s2:one-pool-etr}. These are the precise
novelty claims; the source problem and its algebraic universality are
established results. Since $\NP\subseteq\ETR$, completeness implies that
general pooling threshold decision belongs to $\NP$ only if
$\NP=\ETR$. This conditional obstruction concerns all polynomial-size
certificates, not just certificates listing rational flows. It is
separate from the irrationality and algebraic-degree results proved
below.

\subsection{The algebraic source problem and saturation}\label{s2:source}

An \textsc{ETR-INV} instance consists of variables in $[1/2,2]$ and
equations of the forms
\begin{equation}
 x=1,\qquad x+y=z,\qquad xy=1.\label{s2:etr-inv}
\end{equation}
Variable names may repeat. Abrahamsen, Adamaszek, and Miltzow proved this
problem $\ETR$-complete
\citep[Definition~5 and Theorem~7 of the arXiv version]{abrahamsen2022-the-art-gallery-problem-is}.
Replacing $x=1$ by $x^2=1$ preserves its solutions because all variables
are positive. We henceforth use only addition and inversion equations.
Write $n,m$ for their numbers of variables and equations, and $V(\Phi)$
for the solution set. An extra variable satisfying $u^2=1$ may be added
to ensure $n\ge1$ without changing satisfiability or introducing a choice
in any solution.

The reductions need several prescribed positive totals, although the
threshold instances have zero flow lower bounds. The following elementary
device supplies them. A \emph{marked} input, pool, or product is a node
whose upper capacity is to be attained. Marking is part of the construction,
not an extra type of physical constraint.

\begin{lemma}[Saturation by a profit threshold]\label{s2:saturation}
Let $F_I,F_L,F_J$ be marked nodes with upper capacities $b_u$.
Assign an arc cost of minus one for each of the following groups
containing it: arcs leaving a marked input, arcs leaving a marked pool,
and arcs entering a marked product. With
\[
 \zeta=\sum_{u\in F_I\cup F_L\cup F_J}b_u,
\]
every capacitated flow has profit at most $\zeta$, and it reaches this
threshold exactly when all marked nodes attain their capacities.
Every arc cost lies in $\{0,-1,-2\}$.
\end{lemma}
\begin{proof}
Summing the contributions of each group gives profit
$\sum_{i\in F_I}s_i+\sum_{\ell\in F_L}T_\ell+
\sum_{j\in F_J}t_j$. Each summand is bounded above by its capacity,
so equality in their sum is equivalent to equality in every bound.
A standard arc belongs to at most two groups, including when it is
a bypass. Pool throughput is counted using outgoing arcs only.
\end{proof}

\subsection{One attribute without bypasses}\label{s2:many-pools}

\begin{theorem}\label{s2:etr-complete}
The pooling threshold problem is $\ETR$-complete with all the following
restrictions: one attribute; source qualities in $\{0,1\}$; lower and upper
product quality bounds; no bypasses; zero flow lower bounds; input
out-degree at most two; pool and product in-degree at most two; and
out-degree one at every unmarked pool. Costs lie in $\{0,-1,-2\}$.
Node capacities and quality bounds have numerators and denominators
polynomially bounded in the size of the source instance.
\end{theorem}

We specify the complete construction and prove its correspondence with
$V(\Phi)$. Every node introduced for a use of a gadget is fresh except
for the named connection nodes. Throughout this construction, a
\emph{vacuous} product has quality interval $[0,1]$. A \emph{storage pool}
is a marked pool of capacity $B$, with one quality-one incoming stream,
its own unmarked quality-zero diluent input of capacity $B$, and its own
unmarked vacuous slack product of capacity $B$. The pool has arcs from
both inputs and an arc to its slack product. If its quality-one intake
is $a$, saturation gives diluent intake $B-a$ and quality $a/B$.
The positive integer $B$ will be fixed after counting its other outlets.

For a variable $v$, create a marked quality-one input $s_v$ of capacity
$5/2$, feeding two storage pools $P_v,\bar P_v$. Let $v$ also denote
the flow on $(s_v,P_v)$, and put $\bar v=5/2-v$. At saturation the
other arc carries $\bar v$, and the two pool qualities are $v/B$ and
$\bar v/B$. The notation anticipates, but does not yet assume,
$v\in[1/2,2]$.

\begin{lemma}[Reciprocal emission]\label{s2:emission}
Suppose a storage pool $Q$ has quality-one intake $a$. Add a marked
product $t$ of capacity $2$ and exact quality $1/(2B)$, with inputs
from $Q$ and a new relay pool $H$ of capacity $2$. A marked quality-zero
input $r$ of capacity $2$ feeds $H$ and has exactly one other outgoing
arc, to be connected by the calling gadget. Pool $H$ has only the outlet
to $t$. In any saturated feasible flow,
\begin{equation}
 f_{Qt}=f_{r,\mathrm{other}}=1/a,\qquad
 f_{rH}=f_{Ht}=2-1/a,\qquad a>0.\label{s2:emit-flow}
\end{equation}
In particular, $a\ge1/2$. Conversely, for $a\in[1/2,2]$ these flows
satisfy the new source, relay, and product constraints.
\end{lemma}
\begin{proof}
The quality mass from $H$ is zero, including when its throughput is zero.
The exact quality and saturated total at $t$ give
$(a/B)f_{Qt}=2/(2B)$, so $af_{Qt}=1$.
Its capacity implies $1/a\le2$. Product conservation fixes $f_{Ht}$;
relay conservation fixes $f_{rH}$; saturation of $r$ fixes its remaining
arc. Conversely, $1/a\in[1/2,2]$ makes all four quantities in
\eqref{s2:emit-flow} nonnegative and at most $2$, with the required
totals and quality mass. This checks the gadget conditional on the
specified flow on its connecting arc; the receiving node is checked
by the calling gadget.
\end{proof}

Thus a product's exact quality turns an intake $a$ into a copy of $1/a$
on another input's outgoing arc. We also need to change that arc's source
quality without changing its amount. To convert a delivered quality-zero
amount $d\in[0,2]$ to quality one, route it into a new pool $H_0$ of
capacity $2$, whose only outlet enters a marked vacuous product $t_0$
of capacity $2$. A second pool $H_1$ of capacity $2$ has only an outlet
to $t_0$ and is fed by a new marked quality-one input $r_1$ of capacity
$2$; $r_1$ has one other outgoing arc. Conservation and saturation force
\begin{equation}
 f_{H_0t_0}=d,\qquad f_{r_1H_1}=f_{H_1t_0}=2-d,
 \qquad f_{r_1,\mathrm{other}}=d.\label{s2:flip}
\end{equation}
These nonnegative flows respect all new capacities, and the product is
vacuous. This proves both directions of the conversion. In particular,
it is not a pool--pool arc: all physical arcs still have the standard
input--pool--product form.

Create two more storage pools $R_v,\bar R_v$ for every variable.
Emit from $P_v$ and $\bar P_v$, convert each delivered stream to quality
one by \eqref{s2:flip}, and use these as the quality-one feeds of
$R_v$ and $\bar R_v$. Lemma~\ref{s2:emission} applied to the first
two pools gives $v,\bar v\ge1/2$, hence
\begin{equation}
 v\in[1/2,2],\qquad
 \begin{array}{c|cccc}
 Q&P_v&\bar P_v&R_v&\bar R_v\\ \hline
 Bq_Q&v&\bar v&1/v&1/\bar v\\
 \text{amount delivered by an emission}&1/v&1/\bar v&v&\bar v
 \end{array}\ .\label{s2:four-pools}
\end{equation}
These initial emissions are created even for unused variables. Thus
every variable has the required range independently of its equations.

For an inversion equation $xy=1$, emit from $\bar R_y$ and route its
quality-zero delivery $\bar y$ into a new pool $H_c$ of capacity $5/2$.
Its only outlet enters a marked product $t$ of capacity $5/2$ and exact
quality $2/(5B)$. Add the other product input $(P_x,t)$. Saturation
and conservation give
\[
 f_{H_ct}=\bar y,\qquad f_{P_xt}=5/2-\bar y=y.
\]
The product's quality equation is then
\begin{equation}
 (x/B)y=(5/2)\,2/(5B)=1/B,\label{s2:inversion}
\end{equation}
which is precisely $xy=1$. Conversely, if the equation holds and its
variables lie in $[1/2,2]$, the stated flows are nonnegative and within
the new capacities, and satisfy the exact quality. Nothing changes for
$x=y$.

For an addition $x+y=z$, make separate emissions from $R_x$ and $R_y$,
delivering $x,y$ at quality zero to a new pool $H_+$ of capacity $4$.
Even if $x=y$, the two emissions use distinct relay sources and products.
The only outlet of $H_+$ enters a marked product $t$ of capacity $5/2$
and exact quality $2/(5B)$. Add the second input $(\bar R_z,t)$.
Since $H_+$ carries only quality zero, saturation and the product
quality equation imply
\begin{equation}
 f_{H_+t}=x+y,\qquad
 \frac{5/2-x-y}{B\bar z}=\frac1B,
 \quad\text{hence }x+y=z.\label{s2:addition}
\end{equation}
Conversely, for a solution of the addition equation, these outlet amounts
are $z\le2$ and $\bar z\in[1/2,2]$. They sum to $5/2$, meet its
quality pin, and satisfy the capacity $4$ of $H_+$. Coincidences between
the result variable and a summand cause no ambiguity in this reasoning;
such an equation has no solution in the positive variable range.

It remains to check that many gadget uses can coexist. Count, at each
storage pool, its arcs to products other than its slack product, including
direct arcs in \eqref{s2:inversion} and \eqref{s2:addition}. Let $M$ be
the maximum of these counts and set
\begin{equation}
 B=2M+3.\label{s2:B}
\end{equation}
Initial emissions give $M\ge1$, while $M\le2m+1$. In a solution of
$\Phi$, every one of these non-slack outlets carries at most $2$.
Every quality-one intake in \eqref{s2:four-pools} lies in $[1/2,2]$.
Thus each storage pool can receive its diluent $B-a\in[0,B]$ and
send the residual amount
\begin{equation}
 B-\sum_{t\ne\mathrm{slack}}f_{Qt}\in[3,B]\label{s2:slack}
\end{equation}
to its slack product. Every other capacity was checked locally above.
All product bounds not designated as exact are $[0,1]$; they hold
because every positive stream is a mixture of source qualities zero
and one. At a zero-flow relay choose quality zero or its source quality;
its quality mass is zero either way.

\begin{proof}[Proof of Theorem~\ref{s2:etr-complete}]
Use Lemma~\ref{s2:saturation} with precisely the nodes designated as
marked in the construction.
Any feasible flow meeting its threshold saturates these nodes.
The variable feeds then belong to $[1/2,2]$ by
\eqref{s2:four-pools}, and \eqref{s2:inversion}--\eqref{s2:addition}
recover every equation of $\Phi$. Conversely, a solution of $\Phi$
determines all emissions, conversions, equation gadgets, diluents,
and slack flows by the displayed formulas. The local checks and
\eqref{s2:slack} prove feasibility; saturation gives profit $\zeta$.

There are $O(n+m)$ nodes and arcs. Capacities belong to
$\{2,5/2,4,B\}$, exact quality values are $1/(2B)$ and $2/(5B)$,
and $B\le4m+5$. The threshold is a sum of $O(n+m)$ capacities, hence
has polynomial magnitude and encoding length. Each input has at most
two outlets. Each pool has at most two inputs (the addition pool has
two zero-quality inputs); each product has at most two inputs; each
unmarked pool has one outlet. Costs have the values in
Lemma~\ref{s2:saturation}. Constructing the graph and its rational
data takes polynomial bit time. This proves hardness. Membership is
the quadratic existential formulation in Section~\ref{s1:encoding},
with the linear threshold inequality appended.
\end{proof}

The proof establishes a property stronger than a decision reduction.
\begin{lemma}[Exact correspondence of flows]\label{s2:bijection}
The threshold-feasible arc flows of the construction are in bijection
with $V(\Phi)$. Both maps use rational functions with rational
coefficients whose denominators are nonzero on their domains.
\end{lemma}
\begin{proof}
A solution determines every arc in the construction: variable feeds,
their complements, reciprocal emissions, relay and conversion flows,
equation flows, diluents, and slacks. The preceding proof shows that any
threshold-feasible flow must have exactly these values. The inverse
map reads the arcs $(s_v,P_v)$. The only divisions are by $v$ and
$5/2-v$, both at least $1/2$. These maps are continuous, so they also
give a homeomorphism of the two flow solution sets. The assertion concerns
arc flows, not the irrelevant quality assigned to a zero-flow pool.
\end{proof}

\subsection{Fixed finite data and bounded degrees}\label{s2:finite-data}

\begin{theorem}\label{s2:bounded}
The hardness in Theorem~\ref{s2:etr-complete} persists when every node
has in-degree at most two and out-degree at most three, all capacities
belong to $\{2,5/2,4,7\}$, and all quality bounds belong to
$\{0,1,1/14,2/35\}$. All instance numbers except the profit threshold
belong to fixed finite sets; the threshold is a half-integer of
magnitude $O(n+m)$. The flow correspondence of
Lemma~\ref{s2:bijection} still holds.
\end{theorem}
\begin{proof}
Replace repeated uses of a storage pool by a chain of storage pools,
all of capacity $B=7$. A link consists of an emission, followed by the
quality conversion \eqref{s2:flip}, feeding the next storage pool.
Each new pool has its own capacity-seven diluent and slack product.
Starting with quality-one intake $v$, these links create intakes
\[
 v,\ 1/v,\ v,\ 1/v,\ldots;
 \qquad
 \bar v,\ 1/\bar v,\ \bar v,\ 1/\bar v,\ldots
\]
on two chains rooted at $P_v$ and $\bar P_v$, which share only their
initial variable input. These are chains of logical dependence;
the physical graph remains three-layered and has no pool--pool arcs.

Allocate each equation's requested storage-pool outlet to a distinct
position of the required type on its chain. The requests for an
inversion are one $P_x$-type outlet and one $\bar R_y$-type emission;
those for an addition are two $R$-type emissions and one $\bar R_z$-type
outlet. At a chain position permit at most one equation outlet, besides
its possible link and slack outlet. Process requests sequentially.
If the current position has the wrong type, one new link reaches the
right type. If it already served a request of the right type, two new
links reach a fresh position of that type. Thus each request needs at
most two links. Advance each root at least once, even if its variable
has no equations; this enforces $v,\bar v\ge1/2$ exactly as before.
There are at most $3m$ requests and at most $2n$ additional root links,
so at most $6m+2n$ links suffice.

In a saturated flow, applying Lemma~\ref{s2:emission} and
\eqref{s2:flip} successively along a chain proves the displayed
alternation of intakes. The equation gadgets therefore recover the same
equations, even with repeated variable names. Conversely, for a solution
of $\Phi$, all chain intakes and all non-slack outlets are in $[1/2,2]$.
Each storage pool has at most two non-slack outlets, so its slack is at
least $7-4=3$ and its diluent is at least $7-2=5$. These choices satisfy
every capacity. Each link has constant size, and the rest of the
construction is unchanged, giving $O(n+m)$ nodes and arcs.

For completeness, the local degrees are as follows. Variable, emission,
and conversion inputs have out-degree two; diluents have out-degree
one. Storage pools have in-degree two and out-degree at most three.
Relay and conversion pools have in-degree one and out-degree one;
an addition pool has in-degree two and out-degree one. Every non-slack
product has in-degree two and every slack product has in-degree one.
Substitution of $B=7$ gives the stated finite sets of capacities and
quality bounds. Lemma~\ref{s2:saturation} gives the cost set and the
threshold bound. Every new link is uniquely determined by its initial
intake, so the rational bijection is preserved.
\end{proof}

\begin{corollary}\label{s2:variants}
The preceding threshold classes are $\ETR$-complete also with two
attributes and upper quality bounds only, while keeping binary source
qualities and the same degree restrictions. If exact positive node
throughput contracts are allowed, the corresponding feasibility problems
are $\ETR$-complete. In the bounded-degree feasibility construction all
numerical data belong to fixed finite sets.
\end{corollary}
\begin{proof}
For a source quality $\lambda_i\in\{0,1\}$ add the complementary
attribute $1-\lambda_i$. At a positive-throughput pool the complementary
quality is $1-q_\ell$. At a product the complementary quality mass is
$t_j-m_j$, including at zero flow. Thus a lower bound
$m_j\ge\underline\mu_jt_j$ becomes the upper bound
$t_j-m_j\le(1-\underline\mu_j)t_j$. Zero-flow pools contribute zero
and can receive complementary qualities consistently. No flow or degree
changes, and the set of quality constants remains finite (with their
complements added). For feasibility, impose lower node bounds equal
to the upper capacities precisely at the marked nodes. This implements
the same saturated solution set without an objective or a threshold.
Membership again follows from the quadratic formulation.
\end{proof}

Allowing bypasses or generalized arcs contains these networks as special
cases. In contrast, their zero-lower-bound feasibility problem is trivial
because zero flow is feasible. The threshold construction is not a
feasibility-hardness claim for that zero-lower-bound class. Nor does
complementing an attribute establish $\ETR$-hardness with \emph{one}
upper-bounded attribute. Exact pins supply both directions of
\eqref{s2:inversion} and \eqref{s2:addition}; retaining only an upper
inequality loses that implication. This one-attribute upper-bound-only
$\ETR$ boundary remains unresolved here.

\subsection{One pool with bypasses}\label{s2:one-pool}

\begin{theorem}\label{s2:one-pool-etr}
The pooling threshold problem with a single pool and bypasses is
$\ETR$-complete. Hardness holds with source qualities in $\{0,1\}$,
zero flow lower bounds, input out-degree at most two, product in-degree
at most two, costs in $\{0,-1,-2\}$, and rational numbers of polynomial
magnitude and encoding length. The attribute count and pool degrees
grow with the instance. Lower and upper quality bounds may be replaced
by upper bounds on twice as many attributes. The threshold-feasible
arc flows are again rationally equivalent to the source solution set.
\end{theorem}
\begin{proof}
Start with \textsc{ETR-INV} after replacing $x=1$ by $x^2=1$. Replace
each repeated-summand addition $x+x=z$ by
\[
 xu=1,\qquad ux'=1,\qquad x+x'=z,
\]
using fresh variables $u,x'$. These uniquely equal $1/x,x$, within
$[1/2,2]$, so the replacement preserves solutions by rational maps and
leaves no addition with identical summand names. The construction below
uses this normalized instance.

There is one pool $p$. Each variable $v$ has an unmarked input $s_v$
of capacity $2$ with its only arc to $p$. For each auxiliary index $h$
specified below, add a marked input $a_h$ of capacity $2$, with arcs
to $p$ and to its own marked product $j_h$ of capacity $2$; the only
other input to $j_h$ is $(p,j_h)$. Let $N$ be the total number of
variable and auxiliary inputs and set $B=4N$. Mark the pool with
capacity $B$. Add an unmarked filler input of capacity $B$, whose
only arc goes to $p$, and an unmarked slack product of capacity $B$,
whose only incoming arc is from $p$.

For every non-filler input $s$ introduce an attribute $A_s$, with quality
one at $s$ and zero at every other input. For every addition $x+y=z$
introduce an attribute $A_{x+y}$ supported with quality one on
$s_x,s_y$ only. The normalization makes these binary qualities even
when the original equation had repeated summands. The filler has zero
quality in every attribute. All unspecified product bounds are $[0,1]$.
A \emph{pin} $(j_h,A,\kappa)$ means exact quality $\kappa$ for
attribute $A$ at $j_h$.

An inversion $xy=1$ uses two fresh indices $h_1,h_2$ and the three pins
\begin{equation}
 (j_{h_1},A_{s_x},1/(2B)),\quad
 (j_{h_2},A_{a_{h_1}},1/(2B)),\quad
 (j_{h_2},A_{s_y},1/(2B)).\label{s2:one-inversion-pins}
\end{equation}
An addition $x+y=z$ uses one index $h$ and the two pins
\begin{equation}
 (j_h,A_{s_z},1/B),\qquad(j_h,A_{x+y},1/B).
 \label{s2:one-addition-pins}
\end{equation}
For every variable not appearing in either role of an inversion, add
one more index $h$ and the range pin $(j_h,A_{s_v},1/(2B))$.
This includes unused variables and variables appearing only in additions.

Choose costs by Lemma~\ref{s2:saturation}. If there are $H$ auxiliaries,
the profit threshold is $\zeta=B+4H$. For clarity, the costs on
$(s_v,p)$ and the filler arc are zero, those on $(a_h,p)$ and the
slack outlet are $-1$, and those on $(a_h,j_h)$ and $(p,j_h)$ are $-2$.
In any flow meeting the threshold, put $X_v=f_{s_vp}$ and
$\tau_h=f_{a_hp}$. Source and product saturation give
\begin{equation}
 f_{a_hj_h}=2-\tau_h,\qquad f_{pj_h}=\tau_h.\label{s2:one-copy}
\end{equation}
The pool has throughput $B$, so
$q_p^{A_{s_v}}=X_v/B$, $q_p^{A_{a_h}}=\tau_h/B$, and
$q_p^{A_{x+y}}=(X_x+X_y)/B$. At each pin the bypass input $a_h$
has quality zero in the pinned attribute: in
\eqref{s2:one-inversion-pins} the auxiliary attribute is that of
$a_{h_1}$ at the distinct product $j_{h_2}$. Consequently its pin
equation is
\begin{equation}
 q_p^A\tau_h=2\kappa.\label{s2:pin-product}
\end{equation}

The three inversion pins yield
\[
 X_x\tau_{h_1}=1,\qquad
 \tau_{h_1}\tau_{h_2}=1,\qquad X_y\tau_{h_2}=1.
\]
They force $\tau_{h_1}=X_y$, $\tau_{h_2}=1/X_y$, and $X_xX_y=1$.
Since every variable and auxiliary intake lies in $[0,2]$, these
identities also force $X_x,X_y\in[1/2,2]$, including when $x=y$.
The addition pins give
$X_z\tau_h=(X_x+X_y)\tau_h=2$, so $\tau_h>0$ and
$X_x+X_y=X_z$. Each range pin gives $X_v\tau_h=1$, hence
$X_v\ge1/2$. Thus all variables have the intended range and solve
the source instance.

Conversely, given a source solution $x$, set $X_v=x_v$; set
$\tau_{h_1}=x_y$ and $\tau_{h_2}=1/x_y$ for an inversion;
set $\tau_h=2/x_z$ for an addition; and set $\tau_h=1/x_v$ for
a range pin. In an addition, $x_z=x_x+x_y\in[1,2]$, so
$2/x_z\in[1,2]$. All auxiliary intakes therefore lie in $[1/2,2]$.
Use \eqref{s2:one-copy}, filler intake
$B-\sum_vX_v-\sum_h\tau_h$, and slack outlet
$B-\sum_h\tau_h$. The sum of all non-filler intakes is at most
$2N=B/2$, so filler and slack are both in $[B/2,B]$.
The pool conserves flow and all marked totals are attained.
Equations \eqref{s2:pin-product} hold by substitution. Every unpinned
attribute has average in $[0,1]$, because all source qualities are
binary; its bounds therefore hold, including at the slack product.
This proves threshold feasibility.

Each source solution fixes every flow just listed, and every saturated
flow is forced to have them. Along with the normalization this is the
claimed rational bijection. The graph and attribute count have size
$O(n+m)$; a dense listing of source qualities and product bounds has
size $O((n+m)^2)$, still polynomial. The only nonconstant numerical
values are $B$, $1/B$, $1/(2B)$, and $\zeta$, all with polynomial
encoding length. External degrees follow directly from the arc list.
Membership in $\ETR$ follows from the quadratic formulation. Finally,
with bypasses as well, complementary quality mass is exactly $t_j-m_j$;
the proof of Corollary~\ref{s2:variants} converts the lower bounds to
upper bounds on binary complementary attributes.
\end{proof}

As before, imposing the marked node totals as exact contracts gives
$\ETR$-complete one-pool feasibility. This construction makes no
fixed-attribute or fixed-pool-degree claim. Its many independent
attribute coordinates are essential to the proved scope: the fixed
dimension certificate result below places one pool with a fixed
attribute count in $\NP$.

\subsection{Algebraic degree and the meaning of exactness}\label{s2:degree}

\begin{theorem}\label{s2:algebraic-witness}
For every irrational real algebraic number $\alpha$, there is a rational
pooling instance and a rational profit threshold with exactly one
threshold-feasible arc flow. A designated arc has value
$(\alpha-b)/a$ for rationals $a\ne0,b$, and the field generated by
all its flow values is $\Q(\alpha)$. The instance can be chosen either
with the finite-data and degree restrictions of Theorem~\ref{s2:bounded},
or with one pool and bypasses as in Theorem~\ref{s2:one-pool-etr}.
In particular, rational instances can have only irrational optimal flows
and can require arbitrarily high algebraic degree.
\end{theorem}
\begin{proof}
Choose a nonzero integer polynomial $p$ vanishing at $\alpha$ and
rational $c<\alpha<d$ such that $\alpha$ is its only root in $[c,d]$.
The formula $p(x)=0$, $c\le x\le d$ defines the compact singleton
$\{\alpha\}$. We use the compact \emph{basic closed} case of the
ETR-INV transformations of Abrahamsen and Miltzow
\citep[Lemmas~C--G, underlying Theorem~1]{abrahamsen2019-dynamic-toolbox-for-etrinv}.
Here ``basic closed'' means a conjunction of polynomial equations and
weak inequalities; no elimination of strict inequalities or disjunctions
is needed. To spell out the relevant preservation property, introduce
slacks $s=x-c\ge0$, $t=d-x\ge0$ by equations, and evaluate the
polynomials through addition and multiplication gates. Each new value
is uniquely determined by $x$. Clear rational coefficients to obtain
integer polynomial equations; the added slacks and gates retain a
compact singleton solution set. This is an ETR-COMPACT instance in the
sense required by their Lemma~C, so the compactification step of their
Lemma~B, which need not preserve the solution set, is not used.
Their subsequent transformations scale all these coordinates into a
small interval, shift the computational coordinates near one, replace
multiplication by squaring, and replace squaring by addition and
reciprocal gates. All introduced gate values are uniquely determined
rational functions; the chosen intervals keep their denominators away
from zero. Each reverse transformation recovers the previous equations.
Original coordinates are retained up to rational shifts and nonzero
scalings. Thus the resulting \textsc{ETR-INV} set is a singleton
$\{x^*\}$ with coordinates in $\Q(\alpha)$ and a replacement
coordinate $u$ satisfying $au+b=\alpha$ for rational $a\ne0,b$.
The auxiliary interval promises already hold for all solutions and can
be discarded. In particular, $u=(\alpha-b)/a$ has the same degree as
$\alpha$. Only this basic-closed singleton consequence of the
transformations is used here. It is also the algebraic-field consequence
illustrated by their Corollary~2; the additional step here is its
realization by the restricted physical pooling networks above.

Apply either pooling construction. Its rational bijection produces
exactly one threshold-feasible flow, with all coordinates in
$\Q(\alpha)$, and reads $u$ on a variable feed arc. The chain
refinement retains the initial variable feed; the one-pool
normalization only adds uniquely determined variables. Thus the field
generated by the flows contains $au+b=\alpha$ and is exactly
$\Q(\alpha)$. No subfield of $\R$ omitting $\alpha$ can contain all
flow values. The threshold is an attained global upper bound on profit
by Lemma~\ref{s2:saturation}, so this unique threshold flow is the
unique optimal arc flow. Real algebraic numbers have unbounded degree,
for example the positive roots of $X^d-2$ for integers $d\ge2$, whose
minimal polynomials are irreducible by Eisenstein's criterion at $2$.
\end{proof}

The result concerns required \emph{flow coordinates}; the forcing
objective itself has the rational optimum $\zeta$. Irrational optimal
values already occur in much smaller networks, with one upper quality
bound and every local degree at most two. The following exact example
also illustrates why writing a concentration without its throughput
denominator changes the problem.

\begin{example}[A quadratic irrational optimum]\label{s2:tiny-irrational}
There are inputs $A,B,C$, pools $P,Q$, and products $j_1,j_2$.
Input qualities are $3/4,1/2,0$ and their capacities are $1,2,1$,
respectively. Pool capacities are $2,1$. Each product has capacity one;
its sole quality upper bound is $1/4$ at $j_1$ and $1$ at $j_2$.
All lower flow bounds are zero. The complete arc and cost list is
\[
 \begin{array}{c|rrrrrr}
 \text{arc}&AP&BP&CQ&Pj_2&Pj_1&Qj_1\\ \hline
 \text{cost}&1&2&1&-4&-3&-1
 \end{array}.
\]
Its minimum cost is $-(5+\sqrt3)/2$.
\end{example}
\begin{proof}
Put $a=f_{AP}$, $b=f_{BP}$, $t=a+b$, $z=f_{Pj_1}$, and
$c=f_{Qj_1}=f_{CQ}$. The other pool outlet is $t-z$ and the cost is
\begin{equation}
 a+2b+c-4(t-z)-3z-c=-3a-2b+z.\label{s2:tiny-cost}
\end{equation}
For $t>0$, the quality constraint at $j_1$ is
$(3a+2b)z\le t(z+c)$. Its capacity gives $z+c\le1$,
while the capacity at $j_2$ gives $z\ge\max\{0,t-1\}$.
If $(a,b,z,c)$ is feasible, replacing $z$ by
$z_0=\max\{0,t-1\}$ and $c$ by $1-z_0$ preserves feasibility:
$0\le z_0\le z\le1$, both product capacities hold, the new clean
flow is in $[0,1]$, and
$(3a+2b)z_0\le(3a+2b)z\le t(z+c)\le t$.
It weakly improves \eqref{s2:tiny-cost}. At $t=0$ the mixing flow is
zero and its cost is zero, with no division required.

For $t\le1$, the best profit is at most $3a+2b\le3t\le3$,
attained at $a=1,b=0$. For $t>1$, the normalized choices give profit
$a+t+1$ and quality condition
$(a+2t)(t-1)\le t$. Thus
\[
 0\le a\le\min\{1,h(t)\},\qquad
 h(t)=\frac{t}{t-1}-2t.
\]
Because $1<t\le2$, $a\le1$ ensures $a\le t$, and $b=t-a\le2$;
there are no further intake restrictions. Feasibility is equivalent
to $h(t)\ge0$, or $t\le3/2$. The function $h$ is strictly decreasing
and equals one at $t_0=(1+\sqrt3)/2$. For $1<t\le t_0$ the best
profit is $t+2$, strictly increasing. For $t_0\le t\le3/2$ it is
$2-t+1/(t-1)$, strictly decreasing. The global maximum is therefore
$(5+\sqrt3)/2$, attained with
$a=1$, $b=z=(\sqrt3-1)/2$, and $c=1-b$.
In this family the quality condition at $j_1$ is
$b(3/4+b/2)/(1+b)\le1/4$, equivalently $2b^2+2b-1\le0$.
All displayed flows respect their original capacities.
\end{proof}

The completeness theorem, rather than the irrational example, implies
that general pooling lies in $\NP$ only if $\NP=\ETR$;
the usual containment $\ETR\subseteq\mathsf{PSPACE}$ supplies a
polynomial-space upper bound
\citep[Section~1.1]{abrahamsen2019-dynamic-toolbox-for-etrinv}.
Irrationality alone does not prove failure of $\NP$ membership:
a certificate need not list rational feasible flows.

There is also a precise distinction between exact feasibility and
rational approximate descriptions. For a bounded rational standard
instance with an attained optimum, choose pool qualities in the input
quality box as in Section~\ref{s1:encoding}. All lifted coordinates,
coefficients, and first derivatives of the degree-two constraint
polynomials on a unit enlargement of this box have magnitude at most
$2^{\poly(N)}$, where $N$ is input length. This follows by bounding
each coordinate and coefficient by its rational encoding length and
summing the polynomially many monomials. Consequently a bound
$L=2^{\poly(N)}$ controls the change in every constraint residual and
in the objective when all coordinates change by at most $\delta$.
Round an optimal lifted point to a rational grid of mesh
$\delta\le\epsilon/L$. The rounded description has encoding length
$\poly(N+\log(1/\epsilon))$, its objective differs from the optimum
by at most $\epsilon$, and every equality residual and inequality
violation is at most $\epsilon$. This is an existence statement for
absolute residuals in the specified polynomial formulation, allowing
\emph{all} constraints to have that tolerance. It neither constructs
such a point in polynomial time nor promises an exactly feasible
rational approximation. Rescaling a constraint changes this residual
metric.

\subsection{Linear fibers with a fixed number of parameters}\label{s2:fibers}

We now give the certificate argument used for the ordinary
$\NP$-complete pooling subclasses later in the paper. Its nonlinear
dimension is fixed, but the remaining linear program can have
arbitrarily many coupled variables and rows. The ingredients are
linear programming vertex theory, Cramer's rule, and fixed-dimensional
real-algebraic decision; no novelty is claimed for these ingredients.

\begin{theorem}[A basis-index certificate]\label{s2:fiber-np}
Fix $r$. Consider deciding whether
\begin{equation}
 \eta\in D\subseteq\R^r,\qquad A(\eta)x\le b(\eta),
 \qquad x\in\R^n,\label{s2:fiber-system}
\end{equation}
where $D$ has an explicit rational polynomial description and the
entries of $A,b$ are explicitly represented rational polynomials.
Assume all defining and entry polynomial degrees are polynomially bounded
in input length and
every nonempty fiber in $x$ is bounded. Then feasibility is in $\NP$.
A threshold on an objective linear in $x$, with such polynomial
parameter coefficients, can be included among the rows.
\end{theorem}
\begin{proof}
Suppose first $n>0$ and \eqref{s2:fiber-system} is feasible at
$\eta^*$. Its fiber is a nonempty compact polyhedron and therefore has
a vertex $x^*$. One elementary justification is to minimize its
coordinates successively, retaining each minimizing face; after $n$
steps the retained point is unique, hence extreme. The normals of
active inequalities at that vertex span $\R^n$. Otherwise a nonzero
vector perpendicular to all active normals would give a sufficiently
short feasible segment in both directions, since the finitely many
inactive inequalities have positive slack. This contradicts
extremality, even if the polyhedron is lower-dimensional. Choose $n$
linearly independent active rows.

The certificate consists only of their row indices. Let $C(\eta)$ be
the resulting square coefficient matrix and $d(\eta)$ its right-hand
vector. Form the polynomial determinant and Cramer numerators
\[
 \Delta(\eta)=\det C(\eta),\qquad
 U(\eta)=\operatorname{adj}(C(\eta))d(\eta).
\]
Accept the certificate exactly when the following system in the fixed
$r$ real variables is feasible:
\begin{equation}
 \eta\in D,\qquad \Delta(\eta)^2>0,\qquad
 [A_j(\eta)U(\eta)-b_j(\eta)\Delta(\eta)]\Delta(\eta)\le0
 \quad\text{for every row }j.\label{s2:verifier}
\end{equation}
At $\Delta\ne0$, this is exactly the original row inequality for
$x=U/\Delta$ multiplied by the positive number $\Delta^2$.
It is valid for either determinant sign and excludes the singular
locus. Acceptance therefore produces a feasible point. Conversely,
the selected active basis at $(\eta^*,x^*)$ has nonzero determinant
and satisfies \eqref{s2:verifier}, so some certificate is accepted.

We verify the bit bound as well as the algebra. Let $d_0$ bound the
degrees of entries in $A,b$. Determinants and Cramer numerators have
degree at most $nd_0$, and the row polynomials in
\eqref{s2:verifier} have degree at most $(2n+1)d_0$.
For fixed $r$, dense coefficient arrays at these degrees have
polynomially many entries. Clear input rational denominators first;
their product has polynomial encoding length. If entry coefficients
have magnitude at most $H$ and each entry has at most $M_0$ monomials,
each determinant coefficient is a sum of at most $n!M_0^n$ products
of $n$ coefficients. Its bit length is thus
$O(n\log H+n\log M_0+n\log(n+1))$, polynomial in the input.
The same bound applies to Cramer numerators and, with an additional
polynomial factor, to the verification rows.

These arrays can actually be computed in polynomial time, rather
than merely having polynomial size. Evaluate each determinant on
the tensor grid $\{0,\ldots,nd_0\}^r$, compute its rational
determinants by exact linear algebra, and interpolate one variable
at a time. Both the grid size and the evaluated integer bit lengths
are polynomial. Vandermonde interpolation uses rational linear algebra
on polynomial-size integer matrices with polynomial bit lengths,
so its intermediate and output encoding lengths are polynomial.
The same procedure for the $n$ column-replacement determinants
computes $U$. Constant polynomials cause no difficulty.

Finally, fixed-dimensional real-algebraic decision checks
\eqref{s2:verifier} in polynomial bit time. In particular,
\citet[Theorem~1.3.3 and the bit-size convention in Section~1.3]
{basu1996-on-the-combinatorial-and-algebraic}
give a bound polynomial in the number of polynomials, their degree,
coefficient length, and formula length when the dimension $r$ is fixed.
Strict inequalities such as $\Delta^2>0$ are included in that
sign-condition formulation. The certificate length is
$O(n\log(M+1))$ for $M$ rows. When $n=0$, use the same
fixed-dimensional decision procedure directly, without a basis.
The objective threshold is included before taking a vertex, so
intersecting the original fiber with its halfspace does not affect
boundedness or the argument.
\end{proof}

\begin{corollary}[Pooling certificate scopes]\label{s2:pooling-np}
For standard pooling with rational data and finite flow upper bounds,
feasibility and the linear-cost threshold problem belong to $\NP$
whenever any one of the following quantities is bounded by a fixed
constant:
\begin{enumerate}[label=(\roman*)]
 \item the number of pools times the attribute count;
 \item the number of pools times the affine rank of the input quality
 vectors;
 \item the number of pools times the number of products.
\end{enumerate}
These statements permit arbitrary bypass graphs, flow lower bounds,
exact contracts, and arbitrary rational linear arc costs. They also
give $\NP$ membership for fixed pool and input counts.
\end{corollary}
\begin{proof}
For (i), use all pool qualities as the parameter vector $\eta$.
With these fixed, flow conservation, blending, product quality
constraints, all bounds, and the cost threshold are linear in the
arc flows. Their coefficient polynomials have degree at most one.
Each active pool quality lies in the coordinatewise input-quality
box, and inactive pools may be assigned qualities in that box without
changing feasibility. The parameter domain is therefore a rational
box of fixed dimension, and the flow fibers are bounded.

For (ii), use the rational affine coordinates
$\lambda_i=\lambda_0+Ba_i$ of Section~\ref{s1:encoding} and the
pool coordinates $\eta_\ell$ from \eqref{s1:rank-coordinate}.
At fixed coordinates, that equation and every original product
specification remain linear in flows. The number of parameters is
the pool count times the affine rank, regardless of the number of
original attributes. The affine transformation has polynomial
encoding length; coordinatewise minima and maxima of the $a_i$
bound the active mixtures, and inactive coordinates can be chosen
in the same box. At rank zero this is a linear problem with no
parameters. Empty-input instances have only zero flow and are checked
directly. Since affine rank is at most the input count minus one,
fixed pool and input counts are a special case.

For (iii), parameterize each pool by its outlet fractions
$\theta_{\ell j}\ge0$, with sum one over its allowed outlets.
Remove a pool with no outlet only after checking its pool and incident
arc lower bounds against zero; preserve all external-node constraints.
For a remaining pool let $T_\ell=\sum_i y_{i\ell}$. Substitute
\begin{equation}
 v_{\ell j}=\theta_{\ell j}T_\ell,\qquad
 \text{mass of attribute }k\text{ from }\ell\text{ to }j
   =\theta_{\ell j}\sum_i\lambda_{ik}y_{i\ell}.
 \label{s2:outlet-fractions}
\end{equation}
The remaining variables are intake and bypass flows. With fractions
fixed, input totals, pool totals, product totals, product quality
masses, feed and outlet bounds, and the threshold are all linear
expressions in these variables. The parameter domain is a product
of simplices with at most $|L||J|$ coordinates, and its coefficient
degrees are at most one. A positive-throughput physical pool gives
these fractions by dividing its outlets by $T_\ell$. Conversely,
for a solution of the substituted system with $T_\ell>0$, assign
$q_{\ell k}=\sum_i\lambda_{ik}y_{i\ell}/T_\ell$;
\eqref{s2:outlet-fractions} then gives exactly its physical mass
deliveries. When $T_\ell=0$, all intakes and reconstructed outlets
are zero and any split vector is valid. A pool without outlets was
already checked separately. Finite intake and bypass bounds make
the fibers bounded. Theorem~\ref{s2:fiber-np} proves every case.
\end{proof}

There is one useful certificate scope independent of these products
of dimensions.
\begin{proposition}[One pool without bypasses]\label{s2:no-bypass-np}
One-pool pooling without bypasses belongs to $\NP$ even with
unrestricted input, product, and attribute counts, and with the finite
bounds and linear objectives allowed above.
\end{proposition}
\begin{proof}
Guess a subset $J^+$ of products allowed to receive positive pool flow.
Set all other outlets to zero. Write $T=\sum_i y_i$ and impose the
linear inequalities
\[
 \underline\mu_{jk}T\le\sum_i\lambda_{ik}y_i
   \le\overline\mu_{jk}T\qquad(j\in J^+,\ k=1,\ldots,K),
\]
along with all original linear flow bounds, conservation, and the
threshold. This is a rational linear program decidable in polynomial
bit time. For a physical solution choose $J^+$ to be its positive
outlets; when $T>0$, each such product receives the pool's common
quality, giving exactly these inequalities. If $T=0$, all flows are
zero and the empty choice applies. Conversely, any solution of the
linear program is physical: assign the intake average when $T>0$.
Every positive outlet belongs to $J^+$ and meets its specification;
zero outlets impose no quality condition. At $T=0$, choose an arbitrary
pool quality. This proves soundness even though the guessed set need
not have every one of its outlets positive. The argument is the
active-product linearization in Haugland's proof of Theorem~2
\citep{haugland2016-the-computational-complexity-of-the}, with lower
quality bounds and flow contracts included explicitly.
\end{proof}

These certificate results supply the $\NP$ upper bounds for the
two-pool/two-product hardness construction and the one-pool
fixed-quality bypass constructions developed below. In particular,
one physical quality with lower and upper specifications, or its
two-attribute upper-bound encoding, fits (i). The stronger later
single-upper-quality and constant-physical-data constructions also
fit (i); strong hardness requires their separate reductions and does
not follow from this certificate lemma. None of these results yields
a polynomial algorithm by enumerating bases: there can be
exponentially many candidate row sets. Nor do they give a general
$\NP$ certificate when both the pool count and relevant quality or
outlet dimension grow. The output-fraction proof uses the standard
three-layer model and makes no claim for pool--pool networks.
